% Self-contained snippet: hollowing, absorber optimisation, and homogeneous controls.
% Figure labels assume the manuscript structure proposed in the analysis.

\subsection{Hollowing and absorbing-fill calculations}
\label{sec:hollowing_methods}
Starting from the 120-domain quarter-wave structure with $\lambda_0=470$~nm,
$n_A=1.55$, $n_B=n_m=1$, and material A as the outermost domain, progressive
hollowing was implemented by replacing radial domains from the centre outward by
air while preserving every original radial boundary and the outer particle radius.
Because the B domains already satisfy $n_B=n_m=1$, physically distinct structures
occur after removal of successive A/B pairs. The selected $h=110$ structure leaves
five high-index A domains in the outer reflector. The hollow was subsequently filled
with a homogeneous material $\widetilde n_{\rm fill}=n_{\rm fill}+i\kappa_{\rm fill}$,
with $n_{\rm fill}=1.00$--1.55 and $\kappa_{\rm fill}=3\times10^{-4}$--$2\times10^{-2}$.
Candidate ranking used forward attenuation and backward retention at the same
reference wavelength, blue-band integrals over 400--520~nm, visible-spectrum
correlation, and the CIELAB colour difference $\Delta E_{76}$.

\subsection{Homogeneous effective-medium controls}
\label{sec:homogeneous_controls_methods}
Equal-radius homogeneous spheres were calculated using both a volume-weighted
permittivity and the Bruggeman relation. For the $h=110$ reference, the diameter is
$23.197~\mu$m and the effective indices are $n_{\rm eff}=1.0619$ and 1.0446,
respectively. The homogeneous-sphere extinction coefficient was varied from 0 to 3.
The angle-resolved scattering amplitudes were evaluated on 721 polar angles from
0 to $\pi$. The differential scattering efficiency was normalized so that its
full-sphere integral equalled $Q_{\rm sca}$, and the forward and backward hemisphere
efficiencies were obtained by integration over $0\leq\theta\leq\pi/2$ and
$\pi/2\leq\theta\leq\pi$, respectively. The numerical closure
$Q_F+Q_B-Q_{\rm sca}$ remained below $1.4\times10^{-15}$.

\subsection{Material-efficient hollowing and scattering directionality}
\label{sec:hollowing_results}
The backward structural-colour response was governed predominantly by the outer
radial reflector. Replacing 110 nominal inner domains by air preserved a broad blue
backward-scattering band while leaving only five high-index domains. Weakly
absorbing, intermediate-index hollow fills provided the best compromise between
forward attenuation and colour retention. The leading colour-preserving case,
$n_{\rm fill}=1.20$ and $\kappa_{\rm fill}=5\times10^{-4}$, attenuated forward
scattering at the air-hollow backward peak by 0.55~dB, retained 98.2\% of the
backward response at the same wavelength and 102.0\% of the blue-band backward
integral, and produced $\Delta E_{76}=2.55$. Increasing $\kappa_{\rm fill}$ to
$10^{-3}$ increased the attenuation to 0.67~dB while retaining 95.9\% of the
same-wavelength backward response and 98.2\% of the blue-band integral. A
liquid-like $n_{\rm fill}=1.33$, $\kappa_{\rm fill}=10^{-3}$ fill provided a similar
trade-off with 0.69~dB attenuation and $\Delta E_{76}=4.81$.

The homogeneous controls show why forward scattering cannot be removed by isotropic
absorption alone. Their nonabsorbing response is strongly forward directed, and
moderate extinction reduces the total scattering while increasing the asymmetry
parameter toward unity. Very large extinction subsequently increases scattering
again and reduces $g$, consistent with an opaque-particle diffraction limit.
Consequently, isotropic loss can damp resonant structure but cannot eliminate the
forward diffraction contribution of an optically large spherical particle.

\paragraph{Numerical treatment for plotting.}
For the nominally lossless hollow-onion control, a wavelength was classified as
nonphysical only when $|Q_{\rm abs}|>10^{-6}$ or a plotted quantity was nonfinite.
Only isolated invalid points bounded by valid neighbours were linearly interpolated
in a plotting copy. Original values and an audit table were retained separately.
No amplitude- or gradient-based outlier removal was applied, preventing sharp but
physical Mie resonances from being discarded.

\captionof{figure}{\textbf{Mechanism and homogeneous controls.}
(a) Total scattering spectra of the hollow Bragg onion and equal-radius homogeneous
controls. (b) Forward- and backward-hemisphere efficiencies. (c) Differential
scattering efficiency at 470~nm. (d) Scattering and asymmetry of the Bruggeman
homogeneous control as a function of extinction. Isotropic absorption damps the
total scattering but does not remove the forward diffraction contribution.}
